\documentclass[10pt,letterpaper]{article}

\usepackage[margin=0.58in,top=0.40in,bottom=0.40in]{geometry}
\usepackage{fontspec}
\usepackage[hidelinks]{hyperref}
\usepackage{enumitem}
\usepackage{array}

% Compile with XeLaTeX or LuaLaTeX.
\IfFontExistsTF{Calibri}{\setmainfont{Calibri}}{\setmainfont{Carlito}}

\pagestyle{empty}
\setlength{\parindent}{0pt}
\setlength{\parskip}{0pt}
\setlength{\tabcolsep}{0pt}
\renewcommand{\arraystretch}{1.0}

\newcommand{\bodyfont}{\fontsize{10}{11.1}\selectfont}

\setlist[itemize]{
  leftmargin=1.25em,
  label=\raisebox{0.10ex}{\small\textbullet},
  topsep=2pt,
  itemsep=1pt,
  parsep=0pt,
  partopsep=0pt,
  after=\vspace{4pt}
}

\newcommand{\name}[1]{%
  \begin{center}
    {\fontsize{16.5}{17.5}\selectfont\textbf{#1}}
  \end{center}
  \vspace{-4pt}
}

\newcommand{\contact}[1]{%
  \begin{center}
    {\fontsize{9.6}{10.6}\selectfont #1}
  \end{center}
  \vspace{5pt}
}

% Original-style section divider: line above the section, not under it.
\newcommand{\sectionheader}[1]{%
  \vspace{4pt}
  \noindent\rule{\textwidth}{0.35pt}\par
  \vspace{4pt}
  \noindent{\fontsize{12.5}{13.5}\selectfont\textbf{#1}}\par
  \vspace{3pt}
}

\newcommand{\roleline}[2]{%
  \noindent
  \begin{tabular*}{\textwidth}{@{\extracolsep{\fill}} l r}
    \textbf{#1} & #2 \\
  \end{tabular*}
  \vspace{-2pt}
}

\newcommand{\orgline}[1]{%
  #1\par
  \vspace{1pt}
}

\newcommand{\techline}[1]{%
  #1\par
  \vspace{1pt}
}

\newcommand{\entryspace}{\vspace{2pt}}

\begin{document}
\bodyfont

\name{Akash Karthik}

\contact{San Jose, CA - akashkarthik473@gmail.com - linkedin.com/in/akash-karthik473 - github.com/akashkarthik473 - 9254069507}

\sectionheader{EDUCATION}

\roleline{Bachelor of Science, Computer Science}{Expected May 2027}
\orgline{San Jose State University, San Jose, CA}
\orgline{Relevant Coursework: Data Structures \& Algorithms, Linear Algebra, Discrete Math, Computer Systems, Intro to Circuits}

\sectionheader{TECHNICAL SKILLS}

\techline{\textbf{Languages}: C/C++, Python, Assembly, Java}
\techline{\textbf{Controls}: Fixed-point PI/PID, anti-windup, slew-rate limiting, state estimation, sensor fusion, control-loop tuning}
\techline{\textbf{Embedded}: TTC 60 VCU, STM32H7, PX4, real-time firmware, CAN, SPI, microcontroller I/O}
\techline{\textbf{Tools}: TASKING, TTC Downloader, BusMaster, MoTeC i2 Pro, GDB, Git, CSV telemetry}

\sectionheader{EXPERIENCE}

\roleline{Software Lead}{Jun 2026 - Present}
\orgline{Spartan Racing, Formula SAE\textbar{} San Jose, CA}
\begin{itemize}
  \item Lead software planning for VCU controls, embedded firmware, telemetry, CAN diagnostics, and validation workflows
  \item Coordinate priorities and technical reviews across power limiting, launch control, dashboard, simulation, and vehicle testing
\end{itemize}

\entryspace

\roleline{Software Engineer Intern}{Feb 2026 - Present}
\orgline{Argus Defense \textbar{} San Jose, CA}
\begin{itemize}
  \item Tuned PID flight-control loops on a PX4 flight controller to reduce vertical and horizontal position drift in autonomous flight
  \item Integrated and tuned visual-inertial odometry for onboard state estimation, improving position-hold accuracy
  \item Developed and validated embedded flight-control firmware in C on PX4, iterating control gains against logged flight telemetry
\end{itemize}

\entryspace

\roleline{Software Controls Engineer}{Jul 2025 - Jun 2026}
\orgline{Spartan Racing, Formula SAE\textbar{} San Jose, CA}
\begin{itemize}
  \item Designed a lap-adaptive energy-management algorithm that sets a power cap from cumulative energy versus budget, meeting endurance efficiency targets within \textbf{0.01\%} error
  \item Developed embedded power-limiting firmware with a PI controller on the VCU, maximizing motor output while enforcing FSAE's 80 kW limit
  \item Built CAN diagnostics firmware for live power and torque frames, status bits, and error counters for real-time BusMaster monitoring
\end{itemize}

\entryspace

\roleline{Software Intern}{Aug 2024 - Jun 2025}
\orgline{Spartan Racing, Formula SAE\textbar{} San Jose, CA}
\begin{itemize}
  \item Built deterministic entry and exit state logic for power limiting from measured power, requested torque, and threshold flags
  \item Tuned a PI torque controller to settle with 98\% accuracy and smoothed torque transitions to preserve drivability
  \item \textbf{Results:} 1st in endurance at MIS EV 2025
\end{itemize}

\sectionheader{PROJECT EXPERIENCE}

\roleline{Formula SAE EV Dashboard}{Feb 2026 - Present}
\orgline{Spartan Racing, Formula SAE\textbar{} San Jose, CA}
\techline{C, STM32H7, CAN, SPI}
\begin{itemize}
  \item Built a real-time CAN telemetry display for HV voltage, throttle, cell temperature, and power budget on an LCD
  \item Designed an energy bar and LED alert system giving the driver lap-by-lap feedback on energy budget
\end{itemize}

\entryspace

\roleline{Formula SAE EV Vehicle Simulation}{Jul 2025}
\orgline{Spartan Racing, Formula SAE\textbar{} San Jose, CA}
\techline{Python, Pygame, Matplotlib, PID control, vehicle dynamics, CSV logging}
\begin{itemize}
  \item Built a physics-based FSAE EV simulation to model acceleration using vehicle mass, drag, gearing, and motor torque
  \item Implemented PID-based power limiting to the FSAE 80 kW rule and logged runs to CSV for repeatable tuning and validation
\end{itemize}

\end{document}
